\section{The momentum-cutoff model and its relative bounds}
\label{sec:model}
Set $\hbar=c=1$, $m>0$, $\alpha>0$, $\Lambda=\Lambda_f>0$, and $V=L^3$.
On $\T_L^3=\R^3/L\Z^3$ define
\[
\Gamma_{\Lambda,L}=\{p\in(2\pi/L)\Z^3:|p|\le\Lambda\},\quad
M_L=\#\Gamma_{\Lambda,L},\quad d_L=M_L/V.
\]
The one-particle space $\mathfrak h$ is the span of these plane waves
tensored with $\C^4$, of dimension $D=4M_L$.
Hermitian Dirac matrices $\mathsf a_j,\mathsf b$ obey the Clifford
relations. Set
\[
D_0(p)=\mathsf a\cdot p+\mathsf b m,\quad E(p)=\sqrt{|p|^2+m^2},
\quad P=\mathbf1_{(-\infty,0)}(D_0),\quad \rank P=2M_L.
\]
The finite sea CAR representation $\F=\bigwedge\mathfrak h$ is
unitarily equivalent to the full electron--positron Fock space.
It imposes no Slater restriction and retains all pair terms.
In this representation
\[
H_D=\mathrm d\Gamma(D_0)-\Tr(PD_0)\ge0,\qquad Q=N-2M_L=N_e-N_p.
\]
Electric fermion charge in units of $e$ is $-Q$.
For integral positive nuclear charges $Z_i$, neutrality requires
$Q=Z_{\rm tot}:=\sum_iZ_i$, admissible exactly when $Z_{\rm tot}\le2M_L$.

Let $P_f$ denote the scalar band projection and put
\[
T_k=P_fe^{ikx}P_f\otimes\id_4,\quad
n_P(A)=\mathrm d\Gamma(A)-\Tr(PA),\quad
\rho_k=-n_P(T_k).
\]
Thus $\rho(x)=-:\!\psi^\dagger(x)\psi(x)\!:_P$ and $\rho_0=-Q$.
For $k\ne0$, $\Tr(PT_k)=0$. Only $|k|\le2\Lambda$ can contribute.
The zero-mean periodic Coulomb kernel and form are
\begin{equation}
G_L(x)=\frac1V\sum_{k\ne0}\frac{4\pi}{|k|^2}e^{ikx},\qquad
D_L(f,g)=\frac1V\sum_{k\ne0}\frac{4\pi}{|k|^2}
\overline{\widehat f(k)}\,\widehat g(k).
\label{eq:coulomb}
\end{equation}
The kernel is a distribution, or a heat-summed function off the origin.
Omitting the zero mode and imposing neutrality implements the torus
Gauss constraint. Other charge sectors are auxiliary algebraic sectors,
not isolated nonneutral physical torus systems.

For distinct point nuclei $\nu=\sum_iZ_i\delta_{R_i}$, set
\begin{align}
H_\nu&=H_0+W_\nu+U_{\rm nn},&
W_\nu&=\alpha D_L(\rho,\nu),\\
U_{\rm nn}&=\alpha\sum_{i<j}Z_iZ_jG_L(R_i-R_j),&
H_0&=H_D+\frac{\alpha}{2V}\sum_{k\ne0}v_kW_P(T_{-k},T_k),
\label{eq:model}
\end{align}
where $v_k=4\pi/|k|^2$ and $W_P$ is the full four-field Wick product
relative to the free sea, explicitly defined by \eqref{eq:wick}.
No counterterm is added. The nuclear diagonal is removed by mollifying
each delta, subtracting its \emph{entire periodic Coulomb diagonal},
and taking the limit. This leaves exactly $U_{\rm nn}$.
Keeping a finite self-image term would specify a different energy convention.
The source coupling is the compression of a point Coulomb potential
to the fermion band; the nuclear repulsion itself is not band-truncated.

Each $H_\nu$ is a bounded self-adjoint matrix on all of $\F$ and
commutes with $Q$. These elementary facts do not give uniform stability.
Define
\[
E_\nu(q)=\min\sigma(H_\nu|_{Q=q}),\qquad
E^{\vac}=\min\sigma(H_0|_{Q=0}).
\]
Every comparison uses the same ordering, parameters, and volume.

\begin{theorem}[One nucleus, every box size]\label{thm:one}
Suppose
\begin{equation}
13\alpha\Lambda<\pi^2m.
\label{eq:mass-condition}
\end{equation}
For one nucleus of charge $Z$, every $L>0$ with $Z\le2M_L$, and every $R$,
\begin{equation}
E_{Z\delta_R}(Z)-E^{\vac}\ge-\frac{\alpha Z^2}{2}\kappa_{\Lambda,L}
\ge-\frac{13\alpha\Lambda}{\pi^2}Z^2,\qquad
\kappa_{\Lambda,L}=\frac1V\sum_{0<|k|\le2\Lambda}v_k.
\label{eq:one-bound}
\end{equation}
This is a comparison with the exact interacting vacuum for the full
correlated Fock Hamiltonian.
\end{theorem}

\begin{theorem}[Arbitrary admissible nuclei]\label{thm:many}
Under \eqref{eq:mass-condition}, let $L\ge\Lambda^{-1}$,
$1\le Z_i\le Z_{\max}$, and $Z_{\rm tot}\le2M_L$.
Uniformly in distinct positions and in $K$,
\begin{equation}
E_\nu(Z_{\rm tot})-E^{\vac}\ge
-\alpha\Lambda\left[
24\pi(1+1/\pi)^3Z_{\max}
+\frac{13}{4\sqrt2\,\pi^{3/2}}Z_{\max}^2
\right]K.
\label{eq:many-bound}
\end{equation}
\end{theorem}

At $\alpha=1/137$, the sufficient condition permits
$\Lambda/m<\pi^2\,137/13\simeq104.01$, and $Z_{\max}=26$ is allowed.
The constants and cutoff restriction are not claimed optimal.
A sharper hypothesis is positivity of the effective mass
\eqref{eq:effective-mass} in every box under consideration.
Transverse photons are absent from the model defined so far.

\section{Ordering algebra and the effective mass}
\label{sec:ordering}
\begin{lemma}[Exact Wick identity]\label{lem:wick}
For arbitrary one-particle matrices $A,B$,
\begin{equation}
n_P(A)n_P(B)=W_P(A,B)+n_P(AB-APB-BPA)
+\Tr(PA(1-P)B).
\label{eq:wick}
\end{equation}
Consequently
\begin{align}
H_0&=C+\mathrm d\Gamma(\hf)
+\frac{\alpha}{2V}\sum_{k\ne0}v_k\rho_{-k}\rho_k,
\label{eq:rewrite}\\
\hf&=D_0-\frac\alpha2K_{\rm sign},&
K_{\rm sign}&=\frac1V\sum_{k\ne0}v_kT_{-k}(1-2P)T_k,\\
C&=-\Tr(PD_0)-\frac\alpha2\Ex(P),&
\Ex(P)&=\frac1V\sum_{k\ne0}v_k\Tr(PT_{-k}PT_k).
\label{eq:constant}
\end{align}
This is a rewriting, not a new counterterm.
\end{lemma}
\begin{proof}
Use the CAR contractions
$\langle c_i^\dagger c_j\rangle=P_{ji}$ and
$\langle c_i c_j^\dagger\rangle=\delta_{ij}-P_{ij}$.
After removing the within-bilinear contractions, the single cross
contractions in $c_i^\dagger c_jc_\ell^\dagger c_r$ give
$AB-APB-BPA$, and the double cross contraction is
$\Tr(PA(1-P)B)$. The remaining four-field term is $W_P$.
This proves the polynomial identity \eqref{eq:wick}.

Sum it with $A=T_{-k}$, $B=T_k$ symmetrically in $k$.
Let $K_I=V^{-1}\sum v_kT_{-k}T_k$ and
$K_P=V^{-1}\sum v_kT_{-k}PT_k$.
The single contraction sums to $K_I-2K_P=K_{\rm sign}$.
The double contraction sums to $\Tr(PK_I)-\Ex(P)$.
Expanding $n_P(K_{\rm sign})$ leaves the scalar
\[
\frac\alpha2\Tr(PK_{\rm sign})
-\frac\alpha2[\Tr(PK_I)-\Ex(P)]=-\frac\alpha2\Ex(P).
\]
Adding $H_D$ gives \eqref{eq:rewrite}.
\end{proof}

\begin{lemma}[Uniform mass criterion]\label{lem:mass}
The translation-invariant matrix $\hf$ has the form
\begin{equation}
\hf(p)=\mathsf a\cdot v_{\rm eff}(p)+\mathsf b M_{\rm eff}(p),\qquad
M_{\rm eff}(p)=m-\frac{\alpha}{2V}
\sum_{\substack{r\in\Gamma_{\Lambda,L}\\r\ne p}}
\frac{4\pi}{|p-r|^2}\frac{m}{E(r)}.
\label{eq:effective-mass}
\end{equation}
Its spatial coefficients are Hermitian and scalar in Dirac indices,
and $\Tr\hf=0$. For every $L>0$,
\begin{equation}
\kappa_{\Lambda,L}\le26\Lambda/\pi^2,\qquad
M_{\rm eff}\ge(m-13\alpha\Lambda/\pi^2)\id.
\label{eq:kappa}
\end{equation}
\end{lemma}
\begin{proof}
Insert $1-2P(r)=D_0(r)/E(r)$ in \eqref{eq:rewrite}.
Compression by $T_k$ gives \eqref{eq:effective-mass} and its analogous
vector entry. Every Dirac matrix there is traceless.
Write $k=2\pi n/L$, $R=\Lambda L/\pi$.
The shell $|n|_\infty=j$ contains $24j^2+2$ points, so
\[
\kappa_{\Lambda,L}
=\frac1{\pi L}\sum_{0<|n|\le R}|n|^{-2}
\le\frac{26\lfloor R\rfloor}{\pi L}\le\frac{26\Lambda}{\pi^2}.
\]
For $R<1$ the sum is empty. Use $m/E(r)\le1$ for the mass bound.
\end{proof}

\section{Positivity of the centered fermion trace}
\begin{lemma}[Positive centered trace]\label{lem:positive-trace}
On $\C^n\otimes\C^4$, let
$h=\sum_jv_j\otimes\mathsf a_j+M\otimes\mathsf b$,
where $v_j,M$ are Hermitian and $M>0$.
Consider any finite ordered product $U$ of $e^{-sh}$, $s\ge0$,
and $e^{-i\Phi}$ with $\Phi$ Hermitian and scalar in Dirac indices.
If the total free time is positive, the corresponding centered CAR trace is
\begin{equation}
F=\exp\left(\frac i2\sum_{\rm kicks}\Tr\Phi\right)\det(1+U)>0.
\label{eq:positive-trace}
\end{equation}
The centering is $R(A)=\mathrm d\Gamma(A)-\Tr(A)/2$.
The kicks need not commute with each other or with $h$.
\end{lemma}
\begin{proof}
The Hermitian unitary $\mathsf C_5=\mathsf a_1\mathsf a_2\mathsf a_3
\mathsf b$ anticommutes with all four Dirac matrices and commutes
with scalar kicks. Every factor preserves its indefinite form. Thus
\[
U^\dagger\mathsf C_5U=\mathsf C_5,\qquad
U^\dagger=\mathsf C_5U^{-1}\mathsf C_5.
\]
Put $S=i\sum\Tr\Phi$. Since $\Tr h=0$, $\det U=e^{-S}$ and $S$ is
imaginary. It follows that
\[
\overline{\det(1+U)}=\det(1+U^\dagger)
=\det(U)^{-1}\det(1+U)=e^S\det(1+U).
\]
Therefore $e^{S/2}\det(1+U)$ is real.
The exterior-algebra identities $\Gamma(U_1)\Gamma(U_2)=\Gamma(U_1U_2)$
and $\Tr_{\F}\Gamma(U)=\det(1+U)$ identify it with the centered trace.

It cannot vanish. Along a free interval $f(t)=e^{-th}f(0)$,
\begin{equation}
\frac{\dd}{\dd t}\langle f(t),\mathsf b f(t)\rangle
=-\langle f(t),\{h,\mathsf b\}f(t)\rangle
=-2\langle f(t),M f(t)\rangle<0
\label{eq:beta-monotone}
\end{equation}
for a nonzero vector. Scalar unitary kicks preserve this form.
All factors are invertible. A product with positive total free time
therefore cannot have an eigenvector with an eigenvalue of modulus one,
since its initial and final values of this form would agree.
In particular $-1$ is not an eigenvalue.
Scale all kicks continuously to zero. The real trace never vanishes
and ends at $\det(1+e^{-s_{\rm tot}h})>0$, proving the claim.
\end{proof}

The constant scalar kick $e^{i\theta\id}$ is included in this lemma.
Its centered CAR version is $e^{i\theta(N-D/2)}$.
It permits exact charge projection, which a grand-canonical estimate
by itself would not supply.

\section{Canonical Gaussian domination}
For every spin-scalar compressed multiplication matrix $\Phi$,
$\Tr(P\Phi)=\Tr\Phi/2$: only diagonal momentum entries enter its trace
and $\tr_{\C^4}P(p)=2$.
Choose one $k$ from each pair $\{k,-k\}$ and put
$A_k=(T_k+T_{-k})/2$, $B_k=(T_k-T_{-k})/(2i)$.
Then $\rho_k=-R(A_k)-iR(B_k)$ and
\begin{equation}
\frac{\alpha}{2V}\sum_{k\ne0}v_k\rho_{-k}\rho_k
=\frac\alpha V\sum_{\{k,-k\}}v_k
\left(R(A_k)^2+R(B_k)^2\right).
\label{eq:squares}
\end{equation}
The commutators cancel in the symmetric $k,-k$ sum.

\begin{theorem}[Canonical completed-square comparison]\label{thm:gd}
Assume $M_{\rm eff}>0$. For any real source $\sigma$ with finite
Fourier data on $0<|k|\le2\Lambda$, set
\[
D_{\rm band}(\sigma,\sigma)=V^{-1}
\sum_{0<|k|\le2\Lambda}v_k|\widehat\sigma(k)|^2 .
\]
In every nonempty sector $q$,
\begin{equation}
\min\sigma\!\left[
(H_0+\alpha D_L(\rho,\sigma))|_{Q=q}\right]
\ge E^{\vac}-\frac\alpha2D_{\rm band}(\sigma,\sigma).
\label{eq:gd}
\end{equation}
In particular $E_0(q)\ge E_0(0)$, so the global interacting vacuum
energy is the neutral one.
\end{theorem}
\begin{proof}
Completing \eqref{eq:squares} gives a finite matrix
\[
H(\eta)=C+\mathrm d\Gamma(\hf)
+\frac12\sum_{a=1}^r\lambda_a(R_a+\eta_a)^2,\quad
\lambda_a>0,\quad\eta_a\in\R ,
\]
where $R_a$ is a centered Hermitian spin-scalar bilinear.
There are two components for each $k,-k$ pair, with
$\lambda_a=2\alpha v_k/V$. The scalar source addition is exactly
$\alpha D_{\rm band}(\sigma,\sigma)/2$.
Vanishing compressed modes contribute only a nonnegative scalar square.

Fix $\beta>0$, $\delta=\beta/n$, and use
\[
\mathcal T_n(\eta)=e^{-\beta C}
\left[e^{-\delta\mathrm d\Gamma(\hf)}
\prod_a e^{-\delta\lambda_a(R_a+\eta_a)^2/2}\right]^n.
\]
The spectral theorem gives the exact componentwise identity
\begin{equation}
e^{-\delta\lambda_a(R_a+\eta_a)^2/2}
=\int\frac{e^{-z^2/2}\dd z}{\sqrt{2\pi}}\,
e^{i\sqrt{\delta\lambda_a}zR_a}
e^{i\sqrt{\delta\lambda_a}z\eta_a}.
\label{eq:HS}
\end{equation}
No simultaneous diagonalization of the $R_a$ is used.
Since charge commutes with each factor and is integral,
$P_q=(2\pi)^{-1}\int_0^{2\pi}e^{i\theta(Q-q)}\dd\theta$.
The Gaussian expansions yield
\begin{equation}
\Tr(P_q\mathcal T_n(\eta))
=\frac{e^{-\beta C}}{2\pi}
\int_0^{2\pi}\dd\theta\int\dd\mu(z)\,
e^{-iq\theta}e^{i\ell_\eta(z)}F_n(z,\theta),
\label{eq:canonical}
\end{equation}
where $\ell_\eta$ is real. The trace $F_n$ is a product of centered
free factors and scalar kicks, including $e^{i\theta Q}$.
Lemma~\ref{lem:positive-trace} gives $F_n(z,\theta)>0$ pointwise.
Taking the absolute value therefore gives
\begin{equation}
|\Tr(P_q\mathcal T_n(\eta))|
\le\frac{e^{-\beta C}}{2\pi}\int\dd\theta\,\dd\mu(z)F_n(z,\theta)
=\Tr(P_0\mathcal T_n(0)).
\label{eq:partition}
\end{equation}
The nonsymmetric Trotter operator need not itself be positive.
Finite-matrix norm convergence gives
$0<\Tr_{Q=q}e^{-\beta H(\eta)}
\le\Tr_{Q=0}e^{-\beta H(0)}$ as $n\to\infty$.
Taking $-\beta^{-1}\log$ and then $\beta\to\infty$ proves
\eqref{eq:gd}; the scalar $C$ is identical on both sides.
Setting the source to zero proves the final assertion.
\end{proof}

\begin{corollary}[Interacting dielectric inequality]\label{cor:dielectric}
For real zero-mean $\phi\in H^1(\T_L^3)$ and
$\rho(\phi)=\int\rho(x)\phi(x)\dd x$, every nonempty sector satisfies
\begin{equation}
(H_0+\rho(\phi))|_{Q=q}
\ge\left(E^{\vac}-\frac1{8\pi\alpha}\|\nabla\phi\|_2^2\right)\id.
\label{eq:interacting-dielectric}
\end{equation}
No small-amplitude assumption or assumed interacting gap is needed.
\end{corollary}
\begin{proof}
For nonzero modes take
$\widehat\sigma(k)=|k|^2\widehat\phi(k)/(4\pi\alpha)$.
Then $\rho(\phi)=\alpha D_L(\rho,\sigma)$ and the source cost is
\[
\frac\alpha2D_{\rm band}(\sigma,\sigma)
=\frac1{8\pi\alpha V}\sum_{0<|k|\le2\Lambda}
|k|^2|\widehat\phi(k)|^2
\le\frac1{8\pi\alpha}\|\nabla\phi\|_2^2.
\]
A constant potential couples to $Q$ and is not controlled by this norm.
\end{proof}

\begin{proof}[Proof of Theorem~\ref{thm:one}]
For $\sigma=Z\delta_R$ every Fourier modulus is $Z$ and there is no
nuclear pair term. Apply \eqref{eq:gd} in sector $q=Z$ and use
$D_{\rm band}(\sigma,\sigma)=Z^2\kappa_{\Lambda,L}$
and \eqref{eq:kappa}.
\end{proof}

\section{Point-nucleus electrostatics for arbitrary numbers of nuclei}
\begin{lemma}[Density and heat estimates]\label{lem:heat}
Let $P_t(x)$ denote the positive torus heat kernel,
$G_s=e^{s\Delta}G_L$, and $h_s=G_L-G_s$. Then
\begin{align}
\|\rho(x)\|&\le2d_L,&d_L&\le(\Lambda/\pi+1/L)^3,\label{eq:density}\\
h_s(x)&=4\pi\int_0^s(P_t(x)-V^{-1})\dd t,&
h_s(x)&\ge-4\pi s/V,\quad\|h_s\|_1\le8\pi s,\label{eq:heat}\\
0\le G_s(0)&\le\frac{13}{2\pi^{3/2}\sqrt{s}}.
\label{eq:heat-self}
\end{align}
\end{lemma}
\begin{proof}
The four components $\psi_j(x)/\sqrt{d_L}$ obey four-mode CAR.
Thus $\psi^\dagger(x)\psi(x)$ has spectrum in $[0,4d_L]$ and
free-sea expectation $2d_L$, proving the density estimate.
Counting momenta in a cube proves the bound on $d_L$.
Integrating the heat equation for $G_L$ proves \eqref{eq:heat};
positivity and unit integral of $P_t$ give its two bounds.
With $a=4\pi^2s/L^2$, the maximum-norm shells give
\[
G_s(0)=\frac1{\pi L}\sum_{n\ne0}\frac{e^{-a|n|^2}}{|n|^2}
\le\frac{26}{\pi L}\sum_{j=1}^\infty e^{-aj^2}
\le\frac{26}{\pi L}\int_0^\infty e^{-at^2}\dd t
=\frac{13}{2\pi^{3/2}\sqrt{s}}.
\]
\end{proof}

\begin{proof}[Proof of Theorem~\ref{thm:many}]
Let $\nu_s=e^{s\Delta}\nu$. This is a comparison density; the
Hamiltonian still has point nuclei. Gaussian domination and positivity
of omitted Fourier terms give, in every sector,
\[
H_0+\alpha D_L(\rho,\nu_s)
\ge E^{\vac}-\frac\alpha2D_{\rm band}(\nu_s,\nu_s)
\ge E^{\vac}-\frac\alpha2D_L(\nu_s,\nu_s).
\]
The core error has the operator bound
\begin{equation}
\alpha D_L(\rho,\nu-\nu_s)
=\alpha\sum_i Z_i\int\rho(x)h_s(x-R_i)\dd x
\ge-16\pi\alpha d_LsZ_{\rm tot}.
\label{eq:core}
\end{equation}
This is a norm integral of bounded operators, not a classical
replacement of noncommuting densities.
The remaining nuclear identity is
\begin{align}
U_{\rm nn}-\frac\alpha2D_L(\nu_s,\nu_s)
&=\alpha\sum_{i<j}Z_iZ_jh_{2s}(R_i-R_j)
-\frac\alpha2G_{2s}(0)\sum_iZ_i^2\\
&\ge-\frac{4\pi\alpha s}{V}Z_{\rm tot}^2
-\frac\alpha2G_{2s}(0)\sum_iZ_i^2.
\label{eq:nuclear-heat}
\end{align}
Here $D_L(\nu_s,\nu_s)=\sum_{i,j}Z_iZ_jG_{2s}(R_i-R_j)$.
Neutral capacity $Z_{\rm tot}\le2d_LV$ turns the first error into
$8\pi\alpha d_LsZ_{\rm tot}$. Therefore
\begin{equation}
E_\nu(Z_{\rm tot})-E^{\vac}
\ge-24\pi\alpha d_LsZ_{\rm tot}
-\frac\alpha2G_{2s}(0)\sum_iZ_i^2.
\label{eq:all-s}
\end{equation}
Take $s=\Lambda^{-2}$, use $L\ge\Lambda^{-1}$ and
\eqref{eq:density}--\eqref{eq:heat-self}, then
$Z_{\rm tot}\le Z_{\max}K$ and $\sum_iZ_i^2\le Z_{\max}^2K$.
\end{proof}

The missing ingredient was a vacuum-relative comparison.
Theorem~\ref{thm:gd} supplies it without a volume error, and the
heat argument converts the source self-cost into an $O(K)$ cost.
No operator Graf--Schenker theorem is asserted or needed here.

\subsection{Small boxes remain an actual obstruction}
The one-nucleus theorem holds for all $L$ in the specified diagonal
convention; the many-nucleus theorem cannot.
For $L<2\pi/\Lambda$, only the zero fermion momentum remains.
All nonzero density modes vanish and $H_C=0$.
Two unit nuclei are admissible since $2M_L=2$.
The charge-$2$ state costs $2m$.
There is $x_0\ne0$ on the unit torus with $G_1(x_0)<0$:
otherwise the nonzero mean-zero Green function, positive near the
origin, could not have zero integral.
Choose $R_1-R_2=Lx_0$. Exactly,
\begin{equation}
E_\nu(2)-E^{\vac}=2m+\frac{\alpha}{L}G_1(x_0)
\longrightarrow-\infty\quad(L\downarrow0).
\label{eq:small-box}
\end{equation}
This is a periodic point-source obstruction, not pair collapse.
A fixed positive lower bound on $L$ removes it.
